\documentclass[11pt,a4paper]{article}
\usepackage[utf8]{inputenc}
\usepackage{ctex} % 支持中文
\usepackage{amsmath}
\usepackage{amssymb}
\usepackage{graphicx}
\usepackage{listings}
\usepackage{xcolor}
\usepackage{booktabs}
\usepackage{geometry}
\usepackage{hyperref}

\geometry{left=2.5cm,right=2.5cm,top=3cm,bottom=3cm}

\definecolor{codegreen}{rgb}{0,0.6,0}
\definecolor{codegray}{rgb}{0.5,0.5,0.5}
\definecolor{codepurple}{rgb}{0.58,0,0.82}
\definecolor{backcolour}{rgb}{0.95,0.95,0.92}

\lstset{
    language=C,
    backgroundcolor=\color{backcolour},
    commentstyle=\color{codegreen},
    keywordstyle=\color{magenta},
    numberstyle=\tiny\color{codegray},
    stringstyle=\color{codepurple},
    basicstyle=\ttfamily\small,
    breakatwhitespace=false,
    breaklines=true,
    captionpos=b,
    keepspaces=true,
    numbers=left,
    numbersep=5pt,
    showspaces=false,
    showstringspaces=false,
    showtabs=false,
    tabsize=4,
    extendedchars=false
}

\begin{document}

\section{实验一：并行遗传算法求解旅行商问题 (TSP)}

\subsection{算法背景与分布式并行思路}
本实验针对经典的旅行商问题（TSP），采用分布式并行遗传算法（Island Model GA）进行优化。相比于传统的串行遗传算法，并行方案通过多个进程（岛屿）独立进化，并定期进行优良个体迁移，能够有效维持种群多样性，跳出局部最优解。

\textbf{并行参数设计：}
\begin{itemize}
    \item \textbf{迁移间隔 (Migration Interval)}：每隔一定代数（如 2000 代）进行一次数据交换。
    \item \textbf{迁移数量 (Migration Count)}：每次迁移 2 个优良个体。
    \item \textbf{迁移拓扑 (Migration Topology)}：实现了 \texttt{island} (All-to-All) 和 \texttt{chain} (Ring) 两种模式。
    \item \textbf{迁移策略 (Migration Strategy)}：选择当前种群中最优个体进行迁移（\texttt{best}）。
\end{itemize}

\subsection{核心代码实现与注释}
并行化改进主要体现在 \texttt{main} 函数中的迁移触发和独立的 \texttt{migrate\_if\_needed} 函数。

\begin{lstlisting}[caption={并行 TSP 核心逻辑与注释 (mpitsp.c)}]
// 主循环中的迁移调用
for (GenNum = 0; GenNum < maxGen; GenNum++) {
    for (int i = 0; i < xColony; i++)
        local_improve_one(i); // 执行局部搜索优化个体路径
    select1_local(); // 比较新老个体，保留距离更短的解
    migrate_if_needed(rank, size); // 到达迁移代数时，在进程间交换优良个体
}

// 迁移函数的并行实现
static void migrate_if_needed(int rank, int size) {
    if (GenNum % migration_interval != 0) return;
    
    // 选出最优个体索引并打包
    int send_idx[N_COLONY];
    pick_k_indices(send_idx, migration_count, 1 /* pick_best */);
    
    int send_buf[N_COLONY * CITY];
    for (int m = 0; m < migration_count; m++)
        memcpy(&send_buf[m * xCity], colony[send_idx[m]], xCity * sizeof(int));

    if (topo == TOPO_CHAIN_RING) {
        // 环形拓扑：使用 MPI_Sendrecv 实现相邻进程数据双向交换
        MPI_Sendrecv(send_buf, migration_count * xCity, MPI_INT, next, 100,
                     recv_buf, migration_count * xCity, MPI_INT, prev, 100, 
                     MPI_COMM_WORLD, &st);
    } else {
        // 全连接拓扑：使用 MPI_Allgather 汇总所有进程的优良个体
        MPI_Allgather(send_buf, migration_count * xCity, MPI_INT,
                      all_buf, migration_count * xCity, MPI_INT, MPI_COMM_WORLD);
    }
    // 注入接收到的外部基因，替换本地最差个体
    replace_k_individuals(recv_buf, migration_count, 1 /* replace_worst */);
}
\end{lstlisting}

\subsection{实验结果与统计分析}
实验采用 \texttt{pcb442.tsp} 实例。通过对 \texttt{res/results\_raw.txt} 中的串行（Serial）与并行（MPI, P=4）各 30 次实验数据进行统计，得到如下结果。

\subsubsection{统计计算公式}
1. \textbf{样本均值 ($\bar{x}$)}:
$$\bar{x} = \frac{1}{n} \sum_{i=1}^{n} x_i$$

2. \textbf{样本标准差 ($s$)}:
$$s = \sqrt{\frac{\sum_{i=1}^{n} (x_i - \bar{x})^2}{n-1}}$$

3. \textbf{独立样本 $t$ 检验统计量 ($t$)}:
针对两组样本量相同 ($n_1=n_2=30$) 且方差接近的情况，使用合并方差 $s_p^2$:
$$s_p^2 = \frac{(n_1-1)s_1^2 + (n_2-1)s_2^2}{n_1+n_2-2}$$
$$t = \frac{\bar{x}_1 - \bar{x}_2}{\sqrt{s_p^2 (\frac{1}{n_1} + \frac{1}{n_2})}}$$

\subsubsection{实验数据与计算结果}
根据实验设计，我们在 \texttt{pcb442.tsp} 实例上对串行与并行程序分别运行了 30 次（$n_1=n_2=30$）。详细原始数据记录于 \texttt{res/results\_raw.txt}。

\begin{table}[h]
\centering
\begin{tabular}{@{}lccc@{}}
\toprule
统计项 & 串行 (Serial) & MPI (岛模型) \\ \midrule
样本数 $n$ & 30 & 30 \\
平均值 $\bar{x}$ & 51587.133 & 51351.800 \\
样本方差 $s^2$ & 20837.568 & 17816.924 \\
最好成绩 (Min) & 51241.0 & 51103.0 \\ \bottomrule
\end{tabular}
\caption{TSP 实验统计结果对比表}
\end{table}

\textbf{平均提升分析}：
\begin{itemize}
    \item 均值差值 ($\bar{x}_{serial} - \bar{x}_{mpi}$) = 235.333
    \item 相对解质量提升 $\approx 0.456\%$
\end{itemize}

\subsubsection{Welch $t$ 检验显著性分析}
考虑到两组样本方差不完全相等的情况，本实验采用更为稳健的 \textbf{Welch 双样本 $t$ 检验}（双侧检验）。

\textbf{检验设置}：
\begin{itemize}
    \item 原假设 $H_0$：$\mu_{serial} = \mu_{mpi}$（两组总体均值相同）
    \item 备择假设 $H_1$：$\mu_{serial} \neq \mu_{mpi}$（两组总体均值不同）
\end{itemize}

\textbf{计算过程}：
代入实验统计量到 Welch $t$ 检验公式：
$$t = \frac{\bar{x}_1 - \bar{x}_2}{\sqrt{\frac{s_1^2}{n_1} + \frac{s_2^2}{n_2}}} = 6.556078$$
$$df \approx \frac{\left(\frac{s_1^2}{n_1} + \frac{s_2^2}{n_2}\right)^2}{\frac{(s_1^2/n_1)^2}{n_1-1} + \frac{(s_2^2/n_2)^2}{n_2-1}} = 57.648$$

得到计算结果：
\begin{itemize}
    \item \textbf{$t$-统计量}：6.556078
    \item \textbf{$p$-value}：$1.6646 \times 10^{-8}$
\end{itemize}

\subsubsection{结论分析}
在显著性水平 $\alpha = 0.05$（甚至 $\alpha = 0.01$）下，由于 $p\text{-value} \ll \alpha$，我们**拒绝原假设 $H_0$**。这在统计学上极有力地证明了**并行遗传算法（MPI 岛模型）得到的解质量显著优于串行算法**。

并行化带来的解质量提升主要归功于：
\begin{itemize}
    \item \textbf{多路径搜索}：并行环境下的多个岛屿从不同的随机空间开始探索，显著降低了陷入局部最优的风险。
    \item \textbf{迁移机制的协同效应}：通过 \texttt{Chain} 拓扑定期交换最优个体，使得各个进程能够共享进化路径中的优良基因，从而在有限的迭代次数内更接近全局最优解。
\end{itemize}

\section{实验二：并行 K-means 聚类算法}

\subsection{数据集说明与实际意义}
本次实验使用 \textbf{Iris (鸢尾花) 数据集}。
\begin{itemize}
    \item \textbf{数据特征}：包含 150 个样本，4 个维度（花瓣/花萼的长宽）。
    \item \textbf{实际意义}：鸢尾花数据集是聚类分析的标准 benchmark。在此数据集上运行 K-means 的实际意义在于通过无监督学习，仅凭借植物形态学尺寸，自动将 150 朵花归类为 3 个物种。这对生物分类、自动识别等领域具有研究价值。
\end{itemize}

\subsection{并行化设计思想}
在 K-means 的并行化中，我们面临两种选择：
1. \textbf{任务并行}（每个进程负责一个类中心）：不可取，因为 $K=3$ 无法利用多核优势。
2. \textbf{数据并行}（每个进程负责部分样本点的距离计算）：**本实验采用的方案**。

\textbf{具体实现步骤：}
\begin{itemize}
    \item \textbf{分配阶段}：每个进程仅处理 $1/P$ 规模的数据点。
    \item \textbf{局部统计}：每个进程计算本地样本点到 $K$ 个中心的最近距离，并累加到本地的 \texttt{sum\_local} 和 \texttt{cnt\_local} 数组中。
    \item \textbf{全局同步}：利用 \textbf{\texttt{MPI\_Allreduce}} 函数。将所有进程的局部和进行累加汇总，并同步到所有进程。其操作类似于先求和再广播。
\end{itemize}

\subsection{并行核心代码实现}
\begin{lstlisting}[caption={并行 K-means 数据规约逻辑 (mpikmeans.c)}]
// 1. 各进程计算本地归属并统计部分和
for (int i = begin; i < end; i++) {
    int bestk = find_nearest_centroid(X[i], C);
    cnt_local[bestk]++;
    for (int j = 0; j < DIM; j++) sum_local[bestk][j] += X[i][j];
}

// 2. 使用 MPI_Allreduce 实现全局规约更新
// 将各进程算得的坐标总和 sum_local 累加到 sum_global
MPI_Allreduce(sum_local, sum_global, K*DIM, MPI_DOUBLE, MPI_SUM, MPI_COMM_WORLD);
// 将各进程算得的计数 cnt_local 累加到 cnt_global
MPI_Allreduce(cnt_local, cnt_global, K, MPI_INT, MPI_SUM, MPI_COMM_WORLD);

// 3. 所有进程同步更新中心点
for (int k = 0; k < K; k++) {
    for (int j = 0; j < DIM; j++) C[k][j] = sum_global[k][j] / cnt_global[k];
}
\end{lstlisting}

\subsection{实验结果展示}
通过在本地环境运行串行与并行程序（$P=4$），得到的结果如下：

\subsubsection{串行 K-means 运行结果}
\begin{lstlisting}[backgroundcolor=\color{backcolour}]
Serial k-means on iris: n=149, k=3, dim=4
C0 (count=39): 6.8538 3.0769 5.7154 2.0538
C1 (count=60): 5.8983 2.7483 4.4067 1.4417
C2 (count=50): 5.0060 3.4280 1.4620 0.2460
SSE=76.467679
\end{lstlisting}

\subsubsection{并行 MPI K-means 运行结果}
\begin{lstlisting}[backgroundcolor=\color{backcolour}]
MPI k-means on iris: n=149, k=3, dim=4, P=4
C0 (count=39): 6.8538 3.0769 5.7154 2.0538
C1 (count=60): 5.8983 2.7483 4.4067 1.4417
C2 (count=50): 5.0060 3.4280 1.4620 0.2460
SSE=76.467679
Parallel strategy: data-parallel (each rank handles ~n/P points, Allreduce sums/counts)
\end{lstlisting}

\subsection{结果分析}
从运行结果可以看出：
1. \textbf{数值一致性}：并行版本与串行版本得到的各个簇的中心点坐标、样本计数以及最终的误差平方和 (SSE) 完全一致。这证明了基于 \texttt{MPI\_Allreduce} 的全局规约思路在逻辑上是严谨且正确的。\\
2. \textbf{并行优势}：并行程序通过 \texttt{n/P} 的数据划分，使得每个进程计算距离的开销降低到原来的四分之一。对于像 Iris 这样的小规模数据集，加速效果虽然在毫秒级不明显，但对于百万量级的大数据，这种“数据并行”的架构将展现出显著的性能提升。

\section{实验分析与结论}
1. \textbf{解质量提升}：并行 GA 的“岛屿模型”由于维持了多个独立的进化搜索路径，并通过迁移增强了信息交换，有效避免了串行 GA 易陷入局部最优的问题。T 检验结果明确证明了并行化带来的质量飞跃。\\
2. \textbf{并行效率}：并行 K-means 通过数据均衡划分和高效的规约操作（Allreduce），消除了计算瓶颈，适合大规模聚类任务。
通过本次实验，我们深入理解了如何利用 MPI 的点对点通信（Sendrecv）和集体通信（Allreduce）来改造传统算法，实现高性能计算。

\end{document}

